\setcounter{subsection}{0}
\subsection{Constitucion Española de 1978}

\subsubsection{Introducción}
    Octava constitución de España. Enmarcada en el contexto de la transición española.
    \begin{itemize}
        \item \textbf{1978-10-31} Aprobada en cámaras
        \item \textbf{1978-12-06} Refrendada por el pueblo
        \item \textbf{1978-12-27} Sancionada por el Rey ante las cortes
        \item \textbf{1978-12-29} Publicada en el BOE
    \end{itemize}

    Consta de dos partes, una \textbf{dogmática} donde se recogen los grandes principios y derechos fundamentales, y otra \textbf{orgánica} dónde se recoge la división de poderes, territorial y las competencias.

    Con el clima de la transición en mente, se redactó una constitución rígda para proteger su contenido. Por ello su modificación es costosa. También se intentó un punto medio de compromiso para poder ser aprobada. Por ello es bastante neutra y ambigua. También deja a futuras leyes el detalle de ciertos derechos y deberes.

    La constitución es una ley, y como tal, forma parte del \textbf{ordenamiento jurídico}. También se encuentra por encima de todas las otras normas, en caso de duda, la constitución prevalece. Cabe recurso de inconstitucionalidad contra leyes y disposiciones (art 161.1a) y reforma (art. 166 a 169)

\subsubsection{Principios copnstitucionales}
    Son las normas éticas orientadoras para el conocimiento, interpretación y aplicación de otras normas. Deben inspirar todo el ordenamiento jurídico. 

    El \textbf{Titulo preliminar} define España como un estado \textbf{democrático}, de \textbf{derecho}, \textbf{social}, \textbf{autonómico} y establecido en una \textbf{monarquía parlamentaria}.

    \paragraph{Estado de derecho}
    El Estado de derecho implica una primacía de la ley que regula toda actividad estatal y que se impone a los ciudadanos. Dichas leyes tienen una jerarquía normativa y legal  se elaboran por las cortes. Los ciudadanos son iguales ante la ley

    \paragraph{Estado social}
    Se establece como un estado igualitario. De los españoles ante la ley, de las condiciones que debe promover el estado para que esto sea realizable. El estado debe atender al crecimiento de renta y riqueza y su justa distribución. Para ello garantiza el derecho a la propiedad y a la herencia. Tambien establece la solidaridad entre regiones.

    \paragraph{Estado democrático}
    La soberanía nacional reside en el pueblo, del qe emanan todos los poderes del estado.

    Para ello los partidos políticos, sindicatos, colégios y organizaciones profesionales deben adecuar su estructura y funcionamiento con principios democráticos
    
    \begin{table}[H]
    \caption{Derechos recogidos en la constitución Española}
    \begin{tabular}{>{\bfseries}rll}
    Artículo & Derecho & Protección \\
    \hline
    14 & Igualdad ante la ley & Fundamental \\
    15 & Vida, integridad física y moral & Fundamental \\
    16 & Libertad ideológica & Fundamental \\
    17 & Libertad y seguridad & Fundamental \\
    18 & Honor, intimidad, imágen, domicilio y comunicaciones  & Fundamental \\
    19 & Residiencia y circulación & Fundamental \\
    20 & Expresión e información & Fundamental \\ 
    21 & Reunión pacífica y sin armas & Fundamental \\
    22 & Drecho de asociación & Fundamental \\
    23 & Participación política y cargo público & Fundamental \\
    24 & Tutela judicial & Fundamental \\
    27 & Educación & Fundamental \\
    28 & Sindicación y huelga & Fundamental \\
    29 & Petición & Fundamental \\
    30 & Defender España & Media \\
    31 & contribuir al sostenimiento del gasto público & Media \\
    32 & Contraer matrimonio & Media \\
    33 & Propiedad privada y herencia & Media \\
    34 & Fundación & Media \\
    35 & Trabajo & Media \\
    37 & Negociación colectiva laboral & Media \\
    38 & de Empresa & Media \\
    39 & familia & Mínima \\
    40 & progreso social y económico & Mínima \\
    41 & Seguridad social  & Mínima \\
    42 & Protección en el extranjero & Mínima \\
    43 & a la salud  & Mínima \\
    44 & a la cultura  & Mínima \\
    45 & medio ambiente  & Mínima \\
    46 & Patrimonio histórico, cultural y artístico & Mínima \\
    47 & vivienda  & Mínima \\
    48 & participación de la juventud & Mínima \\
    49 & Atención a las personas con discapacidad & Mínima \\
    50 & Pensiones adecuadas & Mínima \\
    51 & Defensa de los consumidores y usuarios & Mínima \\
    52 & Organizaciones profesionales & Mínima \\
    \end{tabular}
    \end{table}

